\documentclass{article}

\usepackage[english]{babel}
\usepackage[letterpaper, top=2cm, bottom=2cm, left=3cm, right=3cm, marginparwidth=1.75cm]{geometry}

\usepackage{listings}
\lstset{language=Python, basicstyle=\ttfamily\small, frame=single, breaklines=true}

\title{Labwork 2}
\author{Luc Lacotte}

\begin{document}

\maketitle

\section{Introduction}

I am going to play the honest card for this one, I tried looking at the documentation, searching through the source code but I couldn't find anything so some generative AI was used to understand the lab work.

\section{The work}

\subsection{Displaying the graphic card name}

First, I tried doing the commands that were listed in the presentation:

\begin{lstlisting}
print(cuda.detect())
cuda.select_device(0)
\end{lstlisting}

Which displays this:

\begin{lstlisting}
Found 1 CUDA devices
id 0    b'NVIDIA GeForce RTX 3050 Ti Laptop GPU'    [SUPPORTED]
              Compute Capability: 8.6
                   PCI Device ID: 0
                      PCI Bus ID: 1
                        UUID: GPU-6a500019-eb0d-e178-1564-6aee5631d88d
                        Watchdog: Enabled
             FP32/FP64 Performance Ratio: 64
Summary:
    1/1 devices are supported
True
\end{lstlisting}

I am working on my laptop as I have an RTX GPU.

Then I tried looking at the available methods in CUDA. Where I found:

\begin{lstlisting}
get_current_device
\end{lstlisting}

Which gives an object with some attributes from the graphics card and eventually I found how to get only the name with:

\begin{lstlisting}
get_current_device().name
\end{lstlisting}

\subsection{Finding core info}

This is where I started using generative AI because I could not find what I was looking for in the documentation nor the source code.

First, finding the number of multiprocessors was easy, with:

\begin{lstlisting}
device = cuda.get_current_device()
multiprocessor_count = device.MULTIPROCESSOR_COUNT
print(f"Multiprocessor count: {multiprocessor_count}")
>>> Multiprocessor count: 20
\end{lstlisting}

Then I had trouble understanding how I was supposed to get the number of cores because there was not such a thing as \texttt{MULTIPROCESSOR\_COUNT}.

Then I understood that each multiprocessor contained an amount of cores depending on the ``Compute Capability''. This compute capability was given in the \texttt{cuda.detect()} function, 8.6 in my case. With that I have 128 cores per multiprocessor. Thus I needed to multiply both to get the number of cores of my graphics card.

\begin{lstlisting}
device = cuda.get_current_device()
core_count = device.MULTIPROCESSOR_COUNT * 128
print(f"Core count: {core_count}")
>>> Core count: 2560
\end{lstlisting}

\subsection{Finding memory size}

For the memory size, like before I could not find the method I was looking for.

But eventually, there is \texttt{cuda.current\_context().get\_memory\_info()} which gives the current context for memory.

I was then able to get the total amount of memory of my graphics card:

\begin{lstlisting}
_, total = cuda.current_context().get_memory_info()
print(f"Total memory: {total // (1024 ** 2)} MB")
>>> Total memory: 3770 MB
\end{lstlisting}

\section{Full code}
\begin{lstlisting}
from numba import cuda

print(cuda.detect())

cuda.select_device(0)

device = cuda.get_current_device()

# print name
print(device.name)

# multiprocessor count
multiprocessor_count = device.MULTIPROCESSOR_COUNT
print(f"Multiprocessor count: {multiprocessor_count}")

# core count
core_count = device.MULTIPROCESSOR_COUNT * 128 # Compute capability 8.6 => 128 cores per multiprocessor
print(f"Core count: {core_count}")

# memory size
_, total = cuda.current_context().get_memory_info()
print(f"Total memory: {total // (1024 ** 2)} MB")
\end{lstlisting}


\end{document}
